\chapter{Two Services and a Screen}
\label{ch:two-services}

% Act 1. Use cases: UC-01 (spec-hermes#16) and UC-02 (spec-hermes#17), chapter issue spec-hermes#10.
% Code: hermes-credit#1 and #3, hermes-infra#2 (MySQL), hermes-web#1, #2 and #4 (the screens); ADR 0002.
% The chapter is marked by three things: the microservice, MySQL and the problem of indexes, and Angular.
% Already in chapter 1 (refer back, don't repeat): relational vs document in one page; MongoDB's history (2007);
% embedding the documents; Decimal128; the Spring Data repository and derived queries in one box;
% "Is MongoDB better than MySQL?".
% Moved here from chapter 3: Flyway (hermes-credit#5), V1, baseline-on-migrate, ddl-auto: validate.
% Left for chapter 3: dependency injection in depth, the service layer (ScoreCalculator), PUT/PATCH,
% JPA itself, the tests of credit.
% Not here, nobody asked yet: one screen for many services, signals and change detection in depth,
% a static front end on a CDN; embed or reference (16 MB), MongoTemplate, Atlas and DocumentDB.
% Code blocks are copied from hermes-credit at 3a13486 and hermes-web at 7884754.
% Terminal output taken on 2026-11-24: flyway_schema_history from the local hermes-mysql;
% the other examples ran in throwaway containers (mysql:8.4, 100,000 rows; mongo:7, 20,001 documents),
% never in the services. Long output lines are wrapped to fit the page.
% TODO(user): the user's own experience (WM or elsewhere) is not in the chapter yet. Each section has a
% TODO(user) where it would fit. Nothing here claims experience the user did not tell.

In Chapter~\ref{ch:map} you built the first service and chose MongoDB for it.
Now a second team asks for its own tool. The system gets a second service, a
second kind of database, and its first screens. This chapter is about those
three things. At the end you will know why each service keeps its own
database, how a relational database stores and finds data, and how a screen
in the browser talks to a service.

% ===========================================================================
% Mark 1: the microservice
% ===========================================================================
\section{A second team, a second service}

While Sales was getting its tool, the Credit team was looking. Credit is the
team that decides if a farmer gets the money. Here is what they asked for.

\begin{usecase}{UC-02}{Credit analyzes a proposal}{https://github.com/Batchuka/spec-hermes/issues/17}
\textbf{Why now:} analysts score proposals in personal spreadsheets. There is
no history of decisions, and two analysts can decide differently on the same
case.

\textbf{Primary actor:} the credit analyst.

\textbf{Goal:} score a proposal and record an approval or a rejection with its
reason.

\textbf{Preconditions:} the analyst has the proposal in hand, printed or on
the screen of \uc{UC-01}.

\textbf{Main success scenario:}
\begin{enumerate}
  \item The analyst starts a new analysis.
  \item They type the data of the proposal: farmer, income, amount, term.
  \item The system computes a score and a limit.
  \item The analyst approves or rejects, with a reason.
  \item The system stores the decision.
\end{enumerate}

\textbf{Extensions:}
\begin{itemize}
  \item[3a.] Income is missing: the system cannot score and says so.
\end{itemize}

\textbf{Out of scope:} receiving the proposal automatically, because the
analyst retypes it and nobody asked otherwise; and issuing the contract.
\end{usecase}

Look at the first line of the out of scope. The proposal already exists in
the system of Sales, and the analyst types it again by hand. It looks silly.
But nobody asked for the two systems to talk, so they don't. When retyping
starts to hurt, someone will ask, and the system will change then. That is
YAGNI, from Chapter~\ref{ch:map}, doing its job.

\subsection*{Where does the new code go?}

You already have a working service, \texttt{hermes-proposals}. The easy path
is to add a package called \texttt{credit} inside it, with a new controller
and a new table. One application, one deploy, one thing to run.

An application where everything lives together like that is called a
\emph{monolith}, from the Greek for ``one stone''. A monolith is not a bad
word. Many good systems are monoliths, and it is the simplest way to start.

Hermes took the other path: Credit got a new application,
\texttt{hermes-credit}, with its own repository on GitHub. The reason is the
\emph{domain}. A domain is one area of the business, with its own people, its
own words and its own rules. Sales talks about proposals, dealers and
documents. Credit talks about income, scores, limits and decisions. Sales asks
for changes in the proposal; Credit asks for changes in the score. They are
two teams, and they change their minds at different times.

\begin{history}{2014 \textperiodcentered{} ThoughtWorks}
In March 2014, James Lewis and Martin Fowler, both at the consulting company
ThoughtWorks, published an article called ``Microservices'' on Fowler's
website. The word was not theirs. A group of software architects had been
using it since a workshop near Venice in 2011, to name a style they kept
seeing in their clients: one system built as many small applications, each
one running on its own and talking to the others over the network.

The article gave the style a clear description. Two of its points matter in
this chapter. The services are organized around business capabilities, which
is another name for domains. And each service manages its own data.
\end{history}

\begin{concept}{Microservice}
A \emph{microservice} is a small application that serves one domain. It runs
in its own process, has its own database, and talks to the rest of the
system only through its API.

\textbf{What it fixes:} each team changes and deploys its part without waiting
for the others. A bug in Credit does not stop Sales from registering
proposals.

\textbf{What it costs:} more things to run, test and watch. Calls between
services go over the network, which is slower and can fail. And the data of
the system is split across several databases.
\end{concept}

\begin{question}{Isn't it too early for two services?}
Many experienced people would say yes. In 2015 Fowler himself wrote an
article called ``MonolithFirst'', which says you should usually start with a
monolith and split it later.

This book disagrees. Hermes moves fast: Credit asks for its tool right after
Sales. In real projects the second team usually arrives much later, when the
first system is finished and its scope is closed. By then, adding a new
domain inside it means opening a system that was done, and that its team
does not want to touch.

So this book starts separated, one service per domain, and joins things
later if they turn out to be too small to live alone. Joining two small
services is easier than splitting a big one. Divide and conquer.
\end{question}

% TODO(user): your experience with services split by domain (or a monolith that hurt), if any.

\subsection*{Two services on one machine}

The first problem of the second service was not about domains. It was about
ports. A \emph{port} is a number that tells your computer which program
should receive a network request. Spring Boot uses port 8080 unless you say
otherwise. So does \texttt{hermes-proposals}. The two services could not run
at the same time on the same laptop: the second one to start found the port
taken and stopped.

This showed up while the Credit screen was being built, when both services
had to run together for the first time. The fix is one line in the
configuration of Credit:

\begin{code}[language={}]
server:
  port: 8081
\end{code}
\incode{hermes-credit}{900cdb4}{Credit runs on port 8081}

It is a small thing, but it teaches the main lesson of this section: two
services are two separate programs. They don't share memory, ports or
anything else, unless you make them.

\subsection*{Each service owns its data}

Credit needs to store analyses. The quick idea would be to put them in the
MongoDB that proposals already uses. Hermes doesn't, and the reason is a rule.

\begin{concept}{Database per service}
Each service has its own database, and no other service reads or writes it.
If Credit needs something from Sales, it asks the API of Sales. It never
opens the database of Sales.

\textbf{What it fixes:} the shape of the data stays private. Sales can rename
a field, split a collection or change its database, and nothing outside
breaks, as long as its API stays the same. If another service read the
database directly, the database would become a hidden API that nobody
promised to keep.

\textbf{What it costs:} you cannot join the data of two domains in one query.
And some data gets copied: the name of the farmer now lives in both
databases.
\end{concept}

Once each service owns its data, a second question appears: does it have to
be the same \emph{kind} of database? The first decision record of Hermes said
yes, one database server for everyone. Then the data of Credit arrived, and
it had a different shape. A proposal has documents with fields nobody knows
in advance. An analysis has eight fields, always the same. So the decision
changed, and the change was written down in a new record, ADR 0002: each
service picks the database that fits its data. Proposals stays on MongoDB.
Credit uses MySQL.

\begin{center}
\begin{tikzpicture}[x=1cm, y=1cm]
  \node[hf/source, minimum width=26mm] (sales)  at (0,0)    {Sales screen};
  \node[hf/source, minimum width=26mm] (credit) at (0,-2.2) {Credit screen};
  \node[hf/reader, minimum width=30mm] (p) at (4.6,0)    {hermes-proposals\\port 8080};
  \node[hf/reader, minimum width=30mm] (c) at (4.6,-2.2) {hermes-credit\\port 8081};
  \node[hf/store] (mongo) at (9,0)    {MongoDB};
  \node[hf/store] (mysql) at (9,-2.2) {MySQL};
  \draw[hf/ask]  (sales)  -- node[hf/note, above] {HTTP} (p);
  \draw[hf/ask]  (credit) -- node[hf/note, above] {HTTP} (c);
  \draw[hf/save] (p) -- (mongo);
  \draw[hf/save] (c) -- (mysql);
  \draw[hf/pain, dashed] (p.south) -- node[hf/note, right, text=hfAlert] {nothing between them} (c.north);
\end{tikzpicture}
\end{center}

Notice the red dashed line. The two services know nothing about each other.
Each screen talks to one service, and each service talks to one database.

\begin{reflection}{The same farmer is now typed twice, in two databases. Isn't that wrong?}
It is a copy, and copies can disagree. A salesperson can type ``Joao Silva''
and an analyst ``João da Silva''. But the copy is the price of the rule above,
and today it is small: nobody has asked the two systems to agree. When
someone does, the services will need to talk to each other, and that will be
a new request, with its own chapter.
\end{reflection}

% ===========================================================================
% Mark 2: MySQL and the problem of indexes
% ===========================================================================
\section{The relational model}

Chapter~\ref{ch:map} described a relational database in one paragraph: data
in tables, every row with the same columns. Credit is the first service in
Hermes that uses one, so now you will see it properly.

\begin{history}{1970 \textperiodcentered{} San Jose, California}
Before 1970, a program that read a database had to know how the data was
stored on disk: which record pointed to which, and in what order to follow
them. If the storage changed, the programs broke.

Edgar F. Codd, a British mathematician working at the IBM research
laboratory in San Jose, proposed something else. In June 1970 he published
the paper ``A Relational Model of Data for Large Shared Data Banks''. His
idea: show the data to programs only as tables of values, and let the
programs ask for what they want, not for how to find it. The storage could
then change without breaking anyone. IBM was slow to sell it, and the first
commercial relational database came from a small company that later became
Oracle, in 1979. Codd received the Turing Award, the highest prize in
computing, in 1981.
\end{history}

\subsection*{The table of Credit}

In Credit, an analysis is a Java class marked as an \emph{entity}, a class
whose objects are saved as rows of a table. Here it is, without the getters
and setters:

\begin{code}
@Entity
@Table(name = "analyses")
public class Analysis {

	@Id
	@GeneratedValue(strategy = GenerationType.IDENTITY)
	private Long id;

	private String farmer;
	private BigDecimal income;
	private BigDecimal amount;
	private Integer termMonths;
	private Integer score;
	private BigDecimal creditLimit;

	@Enumerated(EnumType.STRING)
	private Decision decision;
	private String reason;
}
\end{code}
\incode{hermes-credit}{53e972e}{the analysis and its table}
\incode{hermes-credit}{93462bb}{the decision and its reason}

The class says what an analysis is in Java. The table says what it is in
MySQL, and it is written in a file of the service, in SQL, the language of
relational databases:

\begin{code}[language=SQL]
CREATE TABLE `analyses` (
  `id` bigint NOT NULL AUTO_INCREMENT,
  `amount` decimal(38,2) DEFAULT NULL,
  `credit_limit` decimal(38,2) DEFAULT NULL,
  `farmer` varchar(255) DEFAULT NULL,
  `income` decimal(38,2) DEFAULT NULL,
  `score` int DEFAULT NULL,
  `term_months` int DEFAULT NULL,
  `decision` enum('APPROVED','REJECTED') DEFAULT NULL,
  `reason` varchar(255) DEFAULT NULL,
  PRIMARY KEY (`id`)
) ENGINE=InnoDB DEFAULT CHARSET=utf8mb4 COLLATE=utf8mb4_0900_ai_ci;
\end{code}
\incode{hermes-credit}{aca9ac9}{the schema as version 1, V1\_\_create\_analyses.sql}

The name of the file, \texttt{V1\_\_create\_analyses.sql}, starts with a
version number. That is on purpose, and Section~\ref{sec:two-shape} tells
the story: the first version of Credit had no such file at all.

Every field of the class became a \emph{column}, with a \emph{type}: whole
numbers are \texttt{int} or \texttt{bigint}, money is \texttt{decimal} with
two digits after the point, text is \texttt{varchar} up to 255 characters,
and the decision can only be one of two words. Every analysis saved becomes
a \emph{row}. The camel case of Java (\texttt{creditLimit}) became the snake
case of SQL (\texttt{credit\_limit}).

The last line names the \emph{primary key}: the column that identifies each
row, so no two rows can have the same \texttt{id}. \texttt{AUTO\_INCREMENT}
means MySQL picks the next number for you. That is why the Java class has
\texttt{@GeneratedValue}.

Here is the same table with two analyses in it. The first one was rejected;
the second one is still waiting for a decision, so its last two cells are
empty:

\begin{center}
\begin{tikzpicture}
  \node[hf/table, font=\sffamily\scriptsize,
        label={[hf/tablename]north west:analyses}] {%
    \begin{tabular}{lllllllll}
      \hfpk{id} & farmer & income & amount & term\_months & score & credit\_limit & decision & reason \\ \hline
      1 & Joao Silva & 50000.00 & 120000.00 & 24 & 41 & 360000.00 & REJECTED & Income too low \\
      2 & Joao & 5000.00 & 10000.00 & 12 & 50 & 18000.00 & & \\
    \end{tabular}};
\end{tikzpicture}
\end{center}

Each column is a fact about the analysis. Each row is one analysis. That
picture is all a relational table is.

One part of the file is still unexplained: \texttt{ENGINE=InnoDB}. MySQL
works in two layers. The top layer understands SQL: it reads your query and
decides how to answer it. The bottom layer, called the \emph{storage engine},
writes the rows to the disk and reads them back. You can pick the engine of
each table, and InnoDB has been the default since MySQL 5.5, in 2010. That
is why you see it at the end of almost every \texttt{CREATE TABLE}. InnoDB is
also the engine that checks foreign keys and keeps a group of changes
all-or-nothing, which is called a \emph{transaction}. The engine that was the
default before it, MyISAM, did neither. Section~\ref{sec:two-further} shows
how InnoDB lays a table out on the disk.

\subsection*{Each fact in one place}

Look at the \texttt{farmer} column. It holds the farmer's name as plain text,
in every analysis. If Joao Silva asks for credit three times, his name is
written three times. If his name was typed wrong, it is wrong in three
places, and fixing one of them does not fix the others.

Codd saw this problem early. In 1970 and 1971 he described a set of rules
called \emph{normal forms}, and the work of changing your tables until they
follow them is called \emph{normalization}. You don't need the rules one by
one to follow this chapter; they are drawn in Section~\ref{sec:two-further},
at the end. The idea behind all of them fits in one sentence: each fact lives
in exactly one place.

The farmer's name is a fact about the farmer, not about the analysis. So the
usual answer is to give the farmers a table of their own. The analysis keeps
only a number that points to the farmer's row.

That number is a \emph{foreign key}: a column that holds the primary key of a
row in another table. The database checks it for you. Here is a normalized
version, run in a throwaway database (it is not the table of Hermes):

\begin{code}[language=SQL]
CREATE TABLE farmers (
  id   bigint NOT NULL AUTO_INCREMENT PRIMARY KEY,
  name varchar(255) NOT NULL
);
CREATE TABLE analyses (
  id        bigint NOT NULL AUTO_INCREMENT PRIMARY KEY,
  farmer_id bigint NOT NULL,
  income    decimal(38,2),
  amount    decimal(38,2),
  FOREIGN KEY (farmer_id) REFERENCES farmers (id)
);
\end{code}

Try to save an analysis for a farmer that does not exist, and MySQL refuses:

\begin{terminal}
mysql> INSERT INTO analyses (farmer_id, income, amount) VALUES (99, 1, 1);
ERROR 1452 (23000): Cannot add or update a child row: a foreign key
constraint fails (`norm`.`analyses`, CONSTRAINT `analyses_ibfk_1`
FOREIGN KEY (`farmer_id`) REFERENCES `farmers` (`id`))
\end{terminal}

To read the analyses with the names, you put the two tables back together
with a \emph{join}. A join matches the rows of two tables by a column they
share:

\begin{terminal}
mysql> SELECT a.id, f.name, a.amount
    -> FROM analyses a
    -> JOIN farmers f ON f.id = a.farmer_id;
id  name        amount
1   Joao Silva  120000.00
2   Joao Silva  450000.00
\end{terminal}

The name is stored once and shown twice. That is the relational model in
one picture: split the data so each fact lives in one place, then join it
back when you read.

\begin{question}{So why doesn't Credit have a farmers table?}
Because nobody asked to manage farmers. \uc{UC-02} says the analyst types the
farmer's name, and that's all. A farmers table would need screens to create,
find and fix farmers, and nobody asked for that. The day someone needs the
history of one farmer, the table will change. Normalization is a tool, not a
duty.
\end{question}

\subsection*{The proposal in tables}

Now you can see the pain that Chapter~\ref{ch:map} only named. Try to store
a proposal of \uc{UC-01} in tables. The proposal is easy: one row with the
farmer, the equipment, the amount and the term. Each document is a row in a
second table, with a foreign key to the proposal.

The fields of the documents are the problem. An invoice has a number and a
value; a proof of income has a year and an income; next month a dealer sends
a type nobody has seen. You can't create columns for fields you don't know.
So you end up with a third table of pairs, one row per field: the document,
the name of the field, and its value as text. To find ``every invoice above
400,000'', you search that table for rows named \texttt{value}, turn the text
into a number, and join three tables. It works, but every query fights the
model. That is why Proposals lives in MongoDB, and Credit, with its eight
fixed fields, lives in MySQL.

% TODO(user): your experience with relational databases (which one, what for), if any.

\section{Who checks the shape}
\label{sec:two-shape}

The two databases of Hermes disagree on one important question: who makes
sure the data has the right shape?

MySQL checks when you write. The table says that \texttt{decision} is
\texttt{APPROVED} or \texttt{REJECTED}. Try something else:

\begin{terminal}
mysql> INSERT INTO analyses (farmer, decision) VALUES ('Joao Silva', 'MAYBE');
ERROR 1265 (01000): Data truncated for column 'decision' at row 1
\end{terminal}

The row is not saved. The message is not very clear (``data truncated''
means the value did not fit the column), but the result is: bad data never
enters the table. This is called \emph{schema on write}. The \emph{schema} is
the description of the shape of the data, and the database applies it every
time you write.

MongoDB, by default, checks nothing. Send it a document with
\texttt{decision: "MAYBE"} and it saves it. The shape is only applied when
some code reads the document and tries to make sense of it. This is called
\emph{schema on read}. You saw it in Chapter~\ref{ch:map}: MongoDB never
received a description of a proposal. The Java record and its Bean
Validation rules are the only schema Proposals has, and they work at the
door of the service, before anything is saved.

\begin{concept}{Schema on write and schema on read}
\textbf{Schema on write:} the database knows the shape and refuses data that
does not fit. \emph{What it fixes:} bad data cannot get in, whatever program
writes it. \emph{What it costs:} every change of shape is a change in the
database, made before the code that uses it.

\textbf{Schema on read:} the database stores what it receives, and the code
that reads decides what it means. \emph{What it fixes:} data whose shape
changes, or is not known in advance, like the documents of \uc{UC-01 3a}.
\emph{What it costs:} the checks live in your code. If one program forgets
them, bad data gets in, and you find out later, when something reads it.
\end{concept}

\begin{question}{Can't MongoDB check the shape too?}
It can. Since version 3.6, in 2017, you can attach a JSON Schema to a
collection, and MongoDB refuses documents that don't match it. Proposals
doesn't use it: the variable fields are the whole point of the service, and
the fixed fields are already checked by the Java code. But it is good to know
the option exists.
\end{question}

\subsection*{The schema, written down}

If MySQL applies the schema, someone has to give the schema to MySQL. A few
pages ago you saw the schema of Credit as a file. In the first version of
the service, that file did not exist.
Nobody wrote a \texttt{CREATE TABLE}. A library called Hibernate, the one
that turns Java objects into rows, read the class \texttt{Analysis} when the
service started and created the table by itself, because the configuration
said \texttt{ddl-auto: update}.

That is fine for the first day, and it goes wrong as soon as the table
changes:

\begin{itemize}
  \item The schema lives nowhere you can read. To see it, you have to ask the
        database.
  \item Changes leave no history. When a column changes, nothing says when,
        why or in what order.
  \item Hibernate guesses. \texttt{update} adds tables and columns that are
        missing, but never removes or renames one, and it changes the
        database every time the service starts.
\end{itemize}

\begin{concept}{Schema migration}
A \emph{schema migration} is one change to the schema, written as a file
with a version number: create a table, add a column, add an index. The files
live with the code of the service and run in order. The schema of the
database is the sum of all the files that ran.

\textbf{What it fixes:} the schema can be read, reviewed and versioned like
code. Every database, on your laptop, in a test or in production, reaches the
same shape by running the same files.

\textbf{What it costs:} a file that already ran must never change. Every
change, even fixing a typo, is a new file.
\end{concept}

\begin{history}{2010 \textperiodcentered{} Flyway}
In 2010, Axel Fontaine, a developer who was tired of changing databases by
hand, released Flyway, a small Java tool that does exactly this. It became
one of the two most used migration tools in the Java world, next to
Liquibase. In 2019 the company Redgate bought it, and its core is still open
source.
\end{history}

In Credit, the migrations live in the folder
\path{src/main/resources/db/migration}. The name of each file follows a rule:
\texttt{V}, the version, two underscores, and a description. The first one
is \texttt{V1\_\_create\_analyses.sql}. When the service starts, Flyway looks
for files that have not run yet on that database, runs them in order, and
writes each one in a table of its own, \path{flyway_schema_history}. It also
keeps a checksum of each file,
a number computed from its content, and refuses to start if a file that
already ran was changed.

The configuration changed in two lines:

\begin{code}[language={}]
spring:
  flyway:
    baseline-on-migrate: true
  jpa:
    hibernate:
      ddl-auto: validate
\end{code}
\incode{hermes-credit}{aca9ac9}{Flyway owns the schema; Hibernate only checks it}

\texttt{ddl-auto: validate} takes the pen away from Hibernate. It no longer
creates or changes anything. When the service starts, it compares the class
with the table and refuses to start if they disagree. Now the schema has one
writer, Flyway, and one checker, Hibernate.

\texttt{baseline-on-migrate} solves a problem that only exists because the
service was already running. The MySQL on the laptop already had the table
\texttt{analyses}, created by Hibernate, and no history table. Flyway refuses
to work on a database like that, because it can't know what is there. The
\emph{baseline} tells it: if you find a database with tables and no history,
write down that version 1 is already applied, and start from there. That is
exactly what the history table of the local database says:

\begin{terminal}
mysql> SELECT version, description, type FROM flyway_schema_history;
version  description            type
1        << Flyway Baseline >>  BASELINE
\end{terminal}

On the laptop, \texttt{V1} never ran; it was marked as done. In the tests,
the MySQL container starts empty, so there \texttt{V1} runs for real, every
time. If the file and the class ever disagree, the tests fail.

From now on, every change to the table of Credit is a new file in that
folder, and the chapter that brings the change shows the file.

\section{Indexes}
\label{sec:two-indexes}

% TODO(user): "the problem of indexes" — your own story goes here first, before the example.

The Credit team asked for a history of decisions. Today the service answers
two questions: ``give me all the analyses'' and ``give me analysis number
42''. Sooner or later an analyst will ask a third one: ``show me every
analysis of this farmer''. Nobody has asked yet, so the service can't do it.
But you can already see what would happen.

To test it, a throwaway MySQL got the same \texttt{analyses} table and
100,000 rows, each one with a different farmer. Then one question, with
\texttt{EXPLAIN ANALYZE} in front of it. \texttt{EXPLAIN ANALYZE} runs the
query and tells you how MySQL found the answer and how long each step took:

\begin{terminal}
mysql> EXPLAIN ANALYZE SELECT * FROM analyses WHERE farmer = 'Farmer 4242';
-> Filter: (analyses.farmer = 'Farmer 4242')  (cost=10116 rows=9947)
   (actual time=1.7..40.3 rows=1 loops=1)
    -> Table scan on analyses  (cost=10116 rows=99473)
       (actual time=0.0793..30.6 rows=100000 loops=1)
\end{terminal}

Read it from the bottom. \emph{Table scan}: MySQL read all 100,000 rows, one
by one, and kept the one that matched. It took 40 milliseconds to find one
row. That sounds fast, but it grows with the table: ten times more analyses,
about ten times slower. And every analyst who opens the history pays it
again.

Now add an index on the column you search by, and ask again:

\begin{terminal}
mysql> CREATE INDEX idx_analyses_farmer ON analyses (farmer);
mysql> EXPLAIN ANALYZE SELECT * FROM analyses WHERE farmer = 'Farmer 4242';
-> Index lookup on analyses using idx_analyses_farmer
   (farmer='Farmer 4242')  (cost=0.35 rows=1)
   (actual time=0.021..0.0221 rows=1 loops=1)
\end{terminal}

\emph{Index lookup}: MySQL went straight to the row. It took 0.02
milliseconds, about two thousand times faster, and it will stay fast when
the table grows.

\begin{concept}{Index}
An \emph{index} is an extra structure that the database keeps next to a
table, sorted by one or more columns, so it can find rows by those columns
without reading the whole table. It works like the index at the back of a
book: you look up the word and go straight to the page.

\textbf{What it fixes:} slow searches. Without an index, finding rows by a
column means reading every row.

\textbf{What it costs:} space on disk, and time on every write. Each
\texttt{INSERT}, \texttt{UPDATE} and \texttt{DELETE} must also update every
index of the table.
\end{concept}

\begin{history}{1970 \textperiodcentered{} Seattle}
The same year Codd published the relational model, Rudolf Bayer and Edward
McCreight, two researchers at Boeing, described a structure for exactly this
job: the \emph{B-tree}. It keeps keys sorted in a short, wide tree, so the
database reaches any key by reading only a few pages of the disk, even with
millions of rows. More than fifty years later, it is still the default.
InnoDB, the storage engine of the table of Credit, uses a variant of it
called the \emph{B+tree}.

It works like a paper dictionary. You never read a dictionary page by page.
You open it, look at the guide words at the top of the page, and jump forward
or back. Each look throws away a big part of the book. A B-tree does the same
with the pages of the disk.

Here is a tiny one, for an index on \texttt{farmer}. Each box is one page.
The top page, the \emph{root}, holds no farmers: Joao and Pedro are only
guide words. Names before Joao go left, names from Joao up to Pedro go to the
middle, and names from Pedro on go right. The pages at the bottom, the
\emph{leaves}, hold the real names, in order, each with the id of its row.

\begin{center}
\begin{tikzpicture}[x=1cm, y=1cm,
    page/.style={hf/table, font=\sffamily\footnotesize},
    edge/.style={hf/flow, draw=black!35},
    path/.style={hf/flow, draw=hfReader, line width=1.4pt}]
  \newcommand\hit[1]{\textbf{\color{hfReader}#1}}
  \node[page] (root) at (0,0) {\begin{tabular}{c|c} Joao & Pedro \end{tabular}};
  \node[page, anchor=north] (l1) at (-4.2,-1.3) {\begin{tabular}{ll} Ana & $\to$ 4 \\ Carlos & $\to$ 6 \end{tabular}};
  \node[page, anchor=north] (l2) at (0,-1.3)    {\begin{tabular}{ll} Joao & $\to$ 1 \\ \hit{Lucas} & \hit{$\to$ 2} \\ Maria & $\to$ 5 \end{tabular}};
  \node[page, anchor=north] (l3) at (4.2,-1.3)  {\begin{tabular}{ll} Pedro & $\to$ 7 \\ Rita & $\to$ 3 \end{tabular}};
  \draw[edge] (root.south west) -- node[hf/note, sloped, above] {before Joao} (l1.north);
  \draw[path] (root.south) -- node[hf/note, right, text=hfReader, pos=0.75] {Joao up to Pedro} (l2.north);
  \draw[edge] (root.south east) -- node[hf/note, sloped, above] {Pedro on} (l3.north);
  \node[hf/note, anchor=west] at (root.east) {\;the root: only guide words};
  \node[hf/note, below=1mm of l3] {the leaves: names and ids};
  \node[hf/table, font=\sffamily\footnotesize] (row) at (0,-3.9)
    {\begin{tabular}{lll} id & farmer & amount \\ \hline 2 & Lucas & 90000.00 \end{tabular}};
  \draw[path] (l2.south) -- (row.north);
  \node[hf/note, anchor=west] at (row.east) {\;the row in the table};
  \node[hf/step, left=2mm of root] {1};
  \node[hf/step, left=2mm of l2] {2};
  \node[hf/step, left=2mm of row] {3};
\end{tikzpicture}
\end{center}

To find Lucas, the database reads three pages:

\begin{enumerate}
  \item \textbf{The root.} Lucas comes after Joao and before Pedro, so it
        follows the middle arrow. The other two pages are never read.
  \item \textbf{The leaf.} Lucas is there, with the id of his row: 2.
  \item \textbf{The row.} The database reads row 2 of the table.
\end{enumerate}

In the drawing each page holds two or three names, so it doesn't look like
much. A real page holds hundreds. With 300 keys in each page, the root points
to 300 pages, those point to 90,000 names, and one more level reaches 27
million. So to find one farmer among the 100,000 analyses of the test above,
the database reads three or four pages, not the 100,000 rows of the table
scan.
\end{history}

You already have one index without asking: the primary key. That is why
``give me analysis number 42'' is fast. MySQL looks up the id in the index of
the primary key, the same way it looked up the farmer above.

\begin{question}{Why not put an index on every column?}
Because each index is paid on every write. A table with ten indexes does ten
extra pieces of work every time you save a row. Index the columns you really
search by, and look at \texttt{EXPLAIN} to check that the database uses them.
In Credit nobody searches by farmer yet, so there is no index on
\texttt{farmer}. When the search arrives, the index arrives with it, as a
new migration: \texttt{V2\_\_index\_analyses\_farmer.sql}.
\end{question}

\subsection*{An index for fields nobody knew}

MongoDB has the same problem, with a twist. The documents of a proposal have
a \texttt{fields} map, and \uc{UC-01 3a} says its names are not known in
advance. A normal index needs the name of the field. How do you index a field
that doesn't exist yet?

MongoDB answers with a \emph{wildcard index}: an index on every field under a
path, whatever its name. In a throwaway MongoDB with 20,001 proposals, a
search for an invoice number examined every document. After the wildcard
index, it examined one. And a field that did not exist when the index was
created, \texttt{anything}, was found through the index too:

\begin{terminal}
> db.proposals.find({ "documents.fields.number": "4242" })
    .explain("executionStats").executionStats.totalDocsExamined
20001
> db.proposals.createIndex({ "documents.fields.$**": 1 })
documents.fields.$**_1
> db.proposals.find({ "documents.fields.number": "4242" })
    .explain("executionStats").executionStats.totalDocsExamined
1
> db.proposals.find({ "documents.fields.anything": "works" })
    .explain("executionStats").executionStats.totalDocsExamined
1
\end{terminal}

The \texttt{\$**} means ``every field below this point''. The cost is the
same as before, only bigger: every field of every document goes into the
index. Proposals doesn't have this index either. Nobody searches inside the
documents yet.

% ===========================================================================
% Mark 3: Angular
% ===========================================================================
\section{A screen for each team}

Until now, the only way to use Hermes was to send JSON to an address, with
a tool like \texttt{curl}. That is fine for a developer and useless for a
salesperson or a credit analyst. With two teams using the system, each one
needed a screen. Chapter~\ref{ch:map} left the screen out on purpose: Spring
Boot and Angular in the same chapter would have been too much at once. The
screens start here.

Both screens live in one Angular project, \texttt{hermes-web}, but they share
nothing except a navigation bar. The Sales screen talks only to
\texttt{hermes-proposals}. The Credit screen talks only to
\texttt{hermes-credit}.

% TODO(user): your experience with Angular or other front ends, if any.

\begin{history}{2009--2016 \textperiodcentered{} Google}
In 2009, Miško Hevery and Adam Abrons started a small library to make web
forms easier to build. Hevery worked at Google, and the library grew there
into AngularJS, released in 2010. It became one of the most used ways to
build screens in the browser.

AngularJS had problems that came from its design, mostly speed in big
screens. So Google did something unusual: instead of fixing it, the team
rewrote it from zero, in TypeScript, and released it in 2016 as Angular. The
two share a name and little else. AngularJS stopped receiving updates at the
end of 2021. When people say ``Angular'' today, they mean the new one. Hermes
uses version 21.
\end{history}

\begin{concept}{Single-page application}
A \emph{single-page application}, or SPA, is a website that loads one page
once and then changes it in the browser, with JavaScript, instead of asking
the server for a new page on every click. It talks to the back end only to
get and send data, as JSON.

\textbf{What it fixes:} the screen feels fast, like a program on your
computer, and the back end only deals with data, not with HTML.

\textbf{What it costs:} a lot of code runs in the browser, and you need a
framework like Angular to keep it organized.
\end{concept}

\subsection*{Components}

An Angular screen is built from \emph{components}. A component is one piece
of the screen: an HTML template, a TypeScript class with its data, and its
CSS. Small components fit together into bigger ones. This is the whole Sales
page:

\begin{code}[language=HTML]
<h1>Register a proposal</h1>
<app-proposal-form (registered)="load()" />
<app-proposal-list [proposals]="proposals()" (selected)="select($event)" />
@if (selected(); as proposal) {
  <app-proposal-detail [proposal]="proposal" />
}
\end{code}
\incode{hermes-web}{2cd8fe6}{the Sales page, one folder per feature}

Three components make the page: a form, a list and a detail. Each tag that
starts with \texttt{app-} is one of them. Two kinds of arrows connect them:

\begin{itemize}
  \item \texttt{[proposals]="..."} is an \emph{input}: the page gives data
        to the list. Square brackets mean ``data goes in''.
  \item \texttt{(registered)="load()"} is an \emph{output}: the form tells
        the page that something happened, and the page reloads the list.
        Round brackets mean ``an event comes out''.
\end{itemize}

\subsection*{Signals: how the list updates by itself}

Look again at \texttt{[proposals]="proposals()"}. Why the parentheses? Because
\texttt{proposals} is not a plain list. It is a \emph{signal}. Here is where
it comes from, in the class of the page:

\begin{code}[language=hfts]
export class ProposalsPage {
  private api = inject(ProposalsApi);

  protected proposals = signal<Proposal[]>([]);

  constructor() {
    this.load();
  }

  protected load() {
    this.api.list().subscribe((list) => this.proposals.set(list));
  }
}
\end{code}
\incode{hermes-web}{2cd8fe6}{the page keeps the list in a signal}

A signal is a box that holds one value and remembers who reads it. You use it
in three ways:

\begin{itemize}
  \item \texttt{signal<Proposal[]>([])} creates the box, with an empty list
        inside.
  \item \texttt{proposals()} reads what is in the box. A signal is read by
        calling it, like a function. That is why the template has
        parentheses.
  \item \texttt{proposals.set(list)} puts a new value in the box.
\end{itemize}

The trick is in the second point. When the template reads
\texttt{proposals()}, Angular takes note: this piece of the screen depends on
that box. Later, when someone calls \texttt{set}, the box tells Angular ``I
changed'', and Angular draws that piece again. Nobody has to tell the list to
refresh. Here is the whole cycle, from the moment the salesperson clicks
Register:

\begin{center}
\begin{tikzpicture}[x=1cm, y=1cm,
    comp/.style={hf/reader, minimum width=30mm, font=\sffamily\footnotesize},
    path/.style={hf/flow, draw=hfReader, line width=1.1pt},
    lbl/.style={hf/note, fill=white, inner sep=1pt}]
  \node[comp] (form) at (0,0)   {\texttt{proposal-form}\\the salesperson\\clicks Register};
  \node[comp] (page) at (5,0) {\texttt{ProposalsPage}\\runs \texttt{load()}};
  \node[hf/reader, minimum width=30mm, font=\sffamily\footnotesize] (api) at (10.8,0)
    {\texttt{hermes-proposals}\\\texttt{GET /proposals}};
  \node[hf/bus, minimum width=30mm, font=\sffamily\footnotesize] (sig) at (5,-2.6)
    {signal \texttt{proposals}\\holds the list};
  \node[comp] (list) at (0,-2.6) {\texttt{proposal-list}\\on the screen,\\with the new row};
  \draw[path] (form) -- node[lbl, above] {\texttt{(registered)}} (page);
  \draw[path] ([yshift=2mm]page.east) -- node[lbl, above] {asks for the list} ([yshift=2mm]api.west);
  \draw[path] ([yshift=-2mm]api.west) -- node[lbl, below] {the list} ([yshift=-2mm]page.east);
  \draw[path] (page) -- node[lbl, right] {\texttt{proposals.set(list)}} (sig);
  \draw[path] (sig) -- node[lbl, above] {``I changed''} (list);
  \node[hf/step, above left=0mm and 0mm of form.north west] {1};
  \node[hf/step, above left=0mm and 0mm of page.north west] {2};
  \node[hf/step, above left=0mm and 0mm of sig.north west] {3};
  \node[hf/step, above left=0mm and 0mm of list.north west] {4};
\end{tikzpicture}
\end{center}

\begin{enumerate}
  \item The salesperson registers a proposal. The form sends the output
        \texttt{(registered)} to the page.
  \item The page runs \texttt{load()}, which asks \texttt{hermes-proposals}
        for the list. The answer comes back a moment later.
  \item \texttt{load()} puts the new list in the signal with
        \texttt{set}.
  \item The signal tells Angular it changed. Angular draws again only what
        reads \texttt{proposals()}: the list. The new proposal appears.
\end{enumerate}

Without the signal, the page would have to find the list component and tell
it to refresh, and every new piece of screen that shows proposals would need
the same call. With the signal, each piece of the screen only says what it
reads, and Angular keeps it up to date. Signals have more to them, but this
is all the two screens of Hermes use.

\subsection*{One class talks HTTP}

None of the components talks to the server. That job belongs to one class
per back end:

\begin{code}[language=hfts]
@Injectable({
  providedIn: 'root',
})
export class ProposalsApi {
  private http = inject(HttpClient);

  list() {
    return this.http.get<Proposal[]>('/proposals');
  }

  get(id: string) {
    return this.http.get<Proposal>(`/proposals/${id}`);
  }

  create(proposal: Omit<Proposal, 'id'>) {
    return this.http.post<Proposal>('/proposals', proposal);
  }
}
\end{code}
\incode{hermes-web}{2cd8fe6}{the only place that talks HTTP to proposals}

In Angular, a class like this is called a \emph{service}. The word is the
same as in ``microservice'', but here it means a class that does one job for
the components. \texttt{HttpClient} is the class of Angular that sends HTTP
requests. Each method returns the request ready to go, and the component
starts it with \texttt{subscribe}, giving the code to run when the answer
arrives.

Look at \texttt{inject(HttpClient)}. The class does not create its
\texttt{HttpClient}; it asks Angular for one. And the components do the same
with \texttt{inject(ProposalsApi)}. This is dependency injection, the same
idea Spring used in Chapter~\ref{ch:map}, when the controller received its
repository in the constructor. \texttt{providedIn: 'root'} tells Angular to
create one \texttt{ProposalsApi} and share it with everyone who asks.

The Credit screen has the same structure, with its own class,
\texttt{AnalysesApi}, that only knows the address of Credit. When the
analyst leaves the income empty, Credit answers 400 with a ProblemDetail, as
in Chapter~\ref{ch:map}, and the form shows its \texttt{detail}: ``Income is
required to score the proposal''. That is \uc{UC-02 3a} on the screen.

\subsection*{Two back ends, one address}

The screen asks for \texttt{/proposals}, not for
\texttt{http://localhost:8080/proposals}. How does the request find the right
service?

During development, Angular runs its own small web server, on port 4200.
That server has a \emph{proxy}: a list of paths, and where to send each one.

\begin{code}[language={}]
{
  "/proposals": {
    "target": "http://localhost:8080",
    "secure": false
  },
  "/analyses": {
    "target": "http://localhost:8081",
    "secure": false
  }
}
\end{code}
\incode{hermes-web}{542e7db}{the dev proxy sends /analyses to Credit}

The browser thinks it talks to one server, the one on port 4200. That
matters, because browsers have a security rule, called \emph{CORS}, that
blocks a page from calling a server at a different address unless that server
allows it. With the proxy, there is only one address, and the services need
no CORS configuration.

There is a side effect. The screens could not live at \texttt{/proposals}
and \texttt{/analyses}, because those paths go to the APIs. So the Sales
screen is \texttt{/sales} and the Credit screen is \texttt{/credit}.

\subsection*{Where am I?}

The first people who used the two screens said they got lost. The page had
two links, ``Sales'' and ``Credit'', but they looked like one word, and
nothing showed which screen was open. Nobody had asked for navigation,
because nobody knew it was needed until they used the screens. It is the
Requirements Uncertainty Principle from Chapter~\ref{ch:map}, in a small
way.

The fix was a navigation bar, and one Angular attribute,
\texttt{routerLinkActive}, which adds a CSS class to the link of the screen
that is open:

\begin{code}[language=HTML]
<nav>
  <a routerLink="/sales" routerLinkActive="active">Sales</a>
  <a routerLink="/credit" routerLinkActive="active">Credit</a>
</nav>
<router-outlet />
\end{code}
\incode{hermes-web}{832878c}{highlight the link of the current screen}

\texttt{routerLink} changes the screen without reloading the page, which is
the single-page application at work. \texttt{<router-outlet />} is the place
where the open screen is drawn.

% ===========================================================================
% Going further: ideas the chapter only names, so the main story is not interrupted.
% ===========================================================================
\section{Going further}
\label{sec:two-further}

This section goes deeper into ideas that the chapter only named, so the main
story would not stop in the middle. You can read it now or come back later.

\subsection*{Normal forms}

A \emph{normal form} is a test that a table passes or fails. Codd defined the
first one in 1970 and the second and third in 1971. Each one includes the one
before it: a table in third normal form is also in the second and the first.

In the drawings below, an \underline{underlined} column is part of the
primary key, a \hfbad{red} value is a fact in the wrong place, and a green
arrow is a foreign key, going from the column to the table it points to.

\textbf{First normal form (1NF): one value per cell.} A \emph{cell} is where
a row meets a column. It must hold one value, never a list. On the left, the
documents of a proposal are written as text in one cell. To find every
proposal with an invoice, you would have to cut that text apart. On the
right, each document is a row of its own.

\begin{center}
\begin{tikzpicture}[x=1cm, y=1cm]
  \node[hf/table, anchor=east, label={[hf/tablename]north west:proposals}] (b) at (3.4,0) {%
    \begin{tabular}{lll}
      \hfpk{id} & farmer & documents \\ \hline
      1 & Joao Silva & \hfbad{invoice, proof of income}
    \end{tabular}};
  \node[hf/note, text=hfAlert, below=3mm of b] {fails 1NF: a list in one cell};
  \draw[hf/flow, draw=black!40] (3.6,0) -- (4.4,0);
  \node[hf/table, anchor=south west, label={[hf/tablename]north west:proposals}] (p) at (4.6,0.35) {%
    \begin{tabular}{ll}
      \hfpk{id} & farmer \\ \hline
      1 & Joao Silva
    \end{tabular}};
  \node[hf/table, anchor=north west, label={[hf/tablename]north west:documents}] (d) at (4.6,-0.35) {%
    \begin{tabular}{lll}
      \hfpk{id} & proposal\_id & type \\ \hline
      10 & 1 & invoice \\
      11 & 1 & proof of income
    \end{tabular}};
  \draw[hf/fk] (d.east) -- ++(0.4,0) |- (p.east);
  \node[hf/note, below=3mm of d] {passes 1NF};
\end{tikzpicture}
\end{center}

MongoDB breaks this rule on purpose: in Chapter~\ref{ch:map} the documents
are a list inside the proposal. That is embedding, and it works there
because the proposal is always read as a whole.

\textbf{Second normal form (2NF): depend on the whole key.} This test only
matters when the primary key has more than one column. On the left, the key
is the pair \texttt{proposal\_id} and \texttt{document\_type}. But
\texttt{farmer} depends only on \texttt{proposal\_id}, half of the key, so it
is repeated in every document of the proposal. On the right, the farmer moves
to a table whose key is the proposal alone.

\begin{center}
\begin{tikzpicture}[x=1cm, y=1cm]
  \node[hf/table, anchor=east, label={[hf/tablename]north west:proposal\_documents}] (b) at (3.4,0) {%
    \begin{tabular}{lll}
      \hfpk{proposal\_id} & \hfpk{document\_type} & farmer \\ \hline
      1 & invoice & \hfbad{Joao Silva} \\
      1 & proof of income & \hfbad{Joao Silva}
    \end{tabular}};
  \node[hf/note, text=hfAlert, below=3mm of b] {fails 2NF: farmer depends on half of the key};
  \draw[hf/flow, draw=black!40] (3.6,0) -- (4.4,0);
  \node[hf/table, anchor=south west, label={[hf/tablename]north west:proposals}] (p) at (4.6,0.35) {%
    \begin{tabular}{ll}
      \hfpk{id} & farmer \\ \hline
      1 & Joao Silva
    \end{tabular}};
  \node[hf/table, anchor=north west, label={[hf/tablename]north west:proposal\_documents}] (d) at (4.6,-0.35) {%
    \begin{tabular}{ll}
      \hfpk{proposal\_id} & \hfpk{document\_type} \\ \hline
      1 & invoice \\
      1 & proof of income
    \end{tabular}};
  \draw[hf/fk] (d.east) -- ++(0.4,0) |- (p.east);
  \node[hf/note, below=3mm of d] {passes 2NF};
\end{tikzpicture}
\end{center}

\textbf{Third normal form (3NF): nothing but the key.} No column may depend
on another column that is not the key. On the left, Credit has added the
farmer's city to the analyses. The city is not a fact about the analysis; it
is a fact about the farmer. When Joao Silva moves, you must change every
analysis of his, or some rows will say one city and some another. On the
right, the farmer and the city live in one row, and each analysis points to
it.

\begin{center}
\begin{tikzpicture}[x=1cm, y=1cm]
  \node[hf/table, anchor=east, label={[hf/tablename]north west:analyses}] (b) at (3.4,0) {%
    \begin{tabular}{llll}
      \hfpk{id} & farmer & city & amount \\ \hline
      1 & Joao Silva & \hfbad{Rio Verde} & 120000.00 \\
      2 & Joao Silva & \hfbad{Rio Verde} & 450000.00
    \end{tabular}};
  \node[hf/note, text=hfAlert, below=3mm of b] {fails 3NF: city depends on the farmer, not on the key};
  \draw[hf/flow, draw=black!40] (3.6,0) -- (4.4,0);
  \node[hf/table, anchor=south west, label={[hf/tablename]north west:farmers}] (p) at (4.6,0.35) {%
    \begin{tabular}{lll}
      \hfpk{id} & name & city \\ \hline
      7 & Joao Silva & Rio Verde
    \end{tabular}};
  \node[hf/table, anchor=north west, label={[hf/tablename]north west:analyses}] (d) at (4.6,-0.35) {%
    \begin{tabular}{lll}
      \hfpk{id} & farmer\_id & amount \\ \hline
      1 & 7 & 120000.00 \\
      2 & 7 & 450000.00
    \end{tabular}};
  \draw[hf/fk] (d.east) -- ++(0.4,0) |- (p.east);
  \node[hf/note, below=3mm of d] {passes 3NF};
\end{tikzpicture}
\end{center}

\begin{concept}{Normal forms}
A short way to remember the three comes from William Kent, in 1983: every
column must give a fact about the key (1NF), the whole key (2NF), and nothing
but the key (3NF).

\textbf{What it fixes:} data that disagrees with itself. When each fact lives
in one place, you change it in one place.

\textbf{What it costs:} more tables, and joins to put them back together
every time you read. That is why some systems copy data on purpose, when
reading fast matters more than writing once. This is called
\emph{denormalization}, and embedding in MongoDB is one kind of it.
\end{concept}

There are more normal forms after the third. Raymond Boyce and Codd defined a
stricter version of the third in 1974, and others came later. In daily work,
a table in third normal form is the usual target.

Now look at the real table of Credit again. Strictly, it passes all three
today: the \texttt{farmer} column holds only a name, and every column says
something about the analysis. But the name is already a fact about the
farmer, copied into every analysis. The first time someone adds a second fact
about the farmer, like the city in the last drawing, the third normal form
breaks.

\subsection*{InnoDB: the table is a tree}

\begin{history}{1995 \textperiodcentered{} Finland}
In 1995, Heikki Tuuri, a Finnish computer scientist, founded a small company
called Innobase and started writing a storage engine: InnoDB. In 2001 it
began to ship with MySQL as an option, and it brought what the default
engine of the time did not have: transactions and foreign keys. In 2005
Oracle bought Innobase. In 2010 Oracle bought Sun Microsystems, which owned
MySQL, and the same year InnoDB became the default engine of MySQL.
\end{history}

In Section~\ref{sec:two-indexes}, the B-tree was an extra structure next to
the table. InnoDB goes one step
further: the table itself is a tree.

It uses the \emph{B+tree}, a variant of the B-tree with two differences. All
the data lives in the \emph{leaves}, the pages at the bottom; the pages
above only hold keys to guide the search. And each leaf is linked to the
next one, so reading a range in order, like every id from 100 to 200, walks
sideways along the bottom without climbing back up.

\textbf{The table is sorted by the primary key.} The leaves of the main tree
hold whole rows, in the order of \texttt{id}. This is called a
\emph{clustered index}: the rows are stored in the order of the key. That is
why a search by id is the fastest search there is. The walk down the tree
ends on the row itself.

\textbf{Every other index is a second tree.} An index on \texttt{farmer} is
a \emph{secondary index}: a separate B+tree, sorted by farmer, whose leaves
hold the farmer and the primary key of the row, not the row. So the search
for Lucas walks two trees. The index on \texttt{farmer} gives the id, 2, and
the main tree turns the id into the row:

\begin{center}
\begin{tikzpicture}[x=1cm, y=1cm,
    page/.style={hf/table, font=\sffamily\scriptsize},
    edge/.style={hf/flow, draw=black!35},
    path/.style={hf/flow, draw=hfReader, line width=1.4pt},
    hit/.style={}]
  \newcommand\hit[1]{\textbf{\color{hfReader}#1}}
  % the secondary index on farmer
  \node[hf/note] at (0,0.75) {index on \texttt{farmer} (secondary)};
  \node[page] (iroot) at (0,0) {\begin{tabular}{c} Lucas \end{tabular}};
  \node[page] (il) at (-1.5,-1.6) {\begin{tabular}{ll} Ana & $\to$ 4 \\ Joao Silva & $\to$ 1 \end{tabular}};
  \node[page] (ir) at (1.5,-1.6)  {\begin{tabular}{ll} \hit{Lucas} & \hit{$\to$ 2} \\ Rita & $\to$ 3 \end{tabular}};
  \draw[edge] (iroot.south) -- (il.north);
  \draw[path] (iroot.south) -- (ir.north);
  % the table itself, clustered by id
  \node[hf/note] at (7.6,0.75) {the table itself, sorted by \texttt{id} (clustered)};
  \node[page] (troot) at (7.6,0) {\begin{tabular}{c} 3 \end{tabular}};
  \node[page] (tl) at (5.5,-1.6) {\begin{tabular}{lll} 1 & Joao Silva & 120000.00 \\ \hit{2} & \hit{Lucas} & \hit{90000.00} \end{tabular}};
  \node[page] (tr) at (9.9,-1.6) {\begin{tabular}{lll} 3 & Rita & 30000.00 \\ 4 & Ana & 75000.00 \end{tabular}};
  \draw[path] (troot.south) -- (tl.north);
  \draw[edge] (troot.south) -- (tr.north);
  % the jump from one tree to the other
  \draw[path] (ir.north east) to[out=60, in=180] node[hf/note, above left, text=hfReader, pos=0.5] {id 2} (troot.west);
  \draw[black!35, densely dotted, line width=0.8pt] (tl.east) -- (tr.west);
  \node[hf/note, below=1mm of tl.south east, anchor=north] {leaves linked in order};
\end{tikzpicture}
\end{center}

Three things follow from this, and they are worth knowing:

\begin{itemize}
  \item A search by a secondary index costs two walks: one in the index, one
        in the table. Still a few pages, but more than a search by id.
  \item The primary key is copied into every secondary index. A short key,
        like the \texttt{bigint} of Credit, keeps every index small.
  \item With \texttt{AUTO\_INCREMENT}, each new row has the biggest id so
        far, so it goes to the end of the tree. A random key, like a UUID,
        lands in the middle of full pages, and InnoDB has to split them to
        make room.
\end{itemize}

% ===========================================================================
\section*{In the interview}

\begin{interview}{Why does each service have its own database?}
Because the database is part of the service, not something shared. In my
project, Hermes, the proposals service and the credit service each own their
data, and no service reads another service's database. If Credit needs
something from Sales, it has to call the API of Sales. The reason is
independence: Sales can change the shape of a proposal, or even change its
database, and nothing outside breaks, as long as the API stays the same. If
Credit read the Mongo collection directly, that collection would become a
hidden contract. Owning the data also let each service pick the right engine.
Proposals has documents with fields nobody knows in advance, so it uses
MongoDB. Credit has analyses with fixed fields, so it uses MySQL. The price
is that I can't join data from two domains in one query, and some data is
copied, like the farmer's name. I wrote that decision down in an ADR.
\end{interview}

\begin{interview}{How do you decide where to put an index?}
I start from the queries, not from the table. An index makes a search fast,
but every insert and update has to maintain it, so I only index columns that
someone really searches by. Then I check with EXPLAIN. Take the credit
table: with a hundred thousand analyses, a search by farmer without an index
is a full table scan. MySQL reads every row, and it takes about forty
milliseconds. With an index on the farmer column, it becomes an index
lookup and takes a few hundredths of a millisecond. The primary key is already
an index, so searching by id is fast from the start. In the real service
there is no index on farmer yet, because nobody searches by farmer yet. When
that request comes, the index comes with it. In MongoDB it is the same idea,
and for fields nobody knows in advance there is the wildcard index.
\end{interview}

\begin{interview}{How does an Angular app talk to a back end?}
Through a service class that uses HttpClient. In my project I have one class
per back end: ProposalsApi talks to the proposals service and AnalysesApi
talks to the credit service. The components never call HTTP themselves. They
get the API class by dependency injection, call a method, and subscribe to
the answer. The answer goes into a signal, and the screen updates by itself.
Errors come back as a ProblemDetail, so the form can show the message of the
server, for example that income is required to score. In development, the
Angular dev server has a proxy: it sends /proposals to one service and
/analyses to the other. The browser sees one address, so I don't need CORS.
The screens are separate, Sales and Credit, with a navigation bar between
them, because each team only works with its own service for now.
\end{interview}
